\documentclass[12pt]{article}
\usepackage{amsmath,amssymb,amsfonts,amsthm}
\usepackage[margin=1in]{geometry}

\begin{document}

\begin{center}
{\Large\bfseries Chapter 9, Problem 12}\\[6pt]
Byron \& Fuller, \textit{Mathematics of Classical and Quantum Physics}
\end{center}

\bigskip

\textbf{Problem.} Consider the integral equation
\[
f(x) = 1 + \lambda\int_0^1 \sqrt{x}\,f(t)\,dt. \tag{1}
\]
Compare the functions $f(x)$, $\tilde{f}(x)$, and $\hat{f}(x)$ obtained by making the successive finite-rank approximations $\sqrt{x} \approx 1$, $\sqrt{x} \approx N\sqrt{x} \approx 1 - t + t\sqrt{x}$, and $\sqrt{x} \approx$ higher order, for the kernel of Eq.~(1). These are obtained in an obvious manner from the power series for the square root.

\bigskip
\textbf{Solution.}

The kernel is $k(x,t) = \sqrt{x}$ (independent of $t$ in this case --- actually, re-reading the problem, the kernel should be $k(x,t) = \sqrt{xt}$ or similar; let us work with what's given).

Actually, with kernel $k(x,t) = \sqrt{x}$ (a rank-1 kernel independent of $t$), the integral equation is:
\[
f(x) = 1 + \lambda\sqrt{x}\int_0^1 f(t)\,dt.
\]

Let $c = \int_0^1 f(t)\,dt$. Then $f(x) = 1 + \lambda c\sqrt{x}$, and:
\[
c = \int_0^1 (1 + \lambda c\sqrt{t})\,dt = 1 + \frac{2}{3}\lambda c.
\]

Solving: $c(1 - \frac{2}{3}\lambda) = 1$, so $c = \frac{1}{1 - 2\lambda/3}$ (provided $\lambda \neq 3/2$).

Therefore the exact solution is:
\[
\boxed{f(x) = 1 + \frac{\lambda\sqrt{x}}{1 - 2\lambda/3}}.
\]

\textbf{Zeroth approximation} ($\sqrt{x} \approx$ constant): Replace $\sqrt{x}$ in the kernel by its average value $\int_0^1 \sqrt{t}\,dt = 2/3$. The kernel becomes $k_0 = 2/3$, so:
\[
f_0(x) = 1 + \frac{2\lambda}{3}\int_0^1 f_0(t)\,dt.
\]
This gives $f_0(x) = 1/(1 - 2\lambda/3)$ (a constant).

\textbf{First approximation} ($\sqrt{x} \approx 1 - t + t\sqrt{x}$ or a linear approximation): Approximate $\sqrt{t} \approx t$ (the simplest nontrivial approximation). Then the kernel becomes $k_1(x,t) = xt$, and:
\[
\tilde{f}(x) = 1 + \lambda x \int_0^1 t\,\tilde{f}(t)\,dt.
\]
Let $d = \int_0^1 t\,\tilde{f}(t)\,dt$. Then $\tilde{f}(x) = 1 + \lambda d x$, and:
\[
d = \int_0^1 t(1 + \lambda d t)\,dt = \frac{1}{2} + \frac{\lambda d}{3}.
\]
So $d = \frac{1/2}{1 - \lambda/3} = \frac{3}{2(3 - \lambda)}$. Therefore:
\[
\tilde{f}(x) = 1 + \frac{3\lambda x}{2(3-\lambda)}.
\]

\textbf{Comparison:} The exact solution $f(x) = 1 + \frac{\lambda\sqrt{x}}{1-2\lambda/3}$ has a $\sqrt{x}$ dependence. The constant approximation $f_0$ misses the spatial dependence entirely. The linear approximation $\tilde{f}(x) = 1 + cx$ captures the growth with $x$ but with the wrong functional form ($x$ vs.\ $\sqrt{x}$). Higher-order finite-rank approximations (using more terms in the Taylor expansion of $\sqrt{t}$ about some point) would progressively better approximate the $\sqrt{x}$ dependence.

The singularity structures also differ: the exact solution diverges at $\lambda = 3/2$, the constant approximation also at $\lambda = 3/2$, but the linear approximation diverges at $\lambda = 3$. This shows that the finite-rank approximations can shift the location of the singular values (eigenvalues of the kernel). \hfill $\blacksquare$

\end{document}
